\section{Evaluated Approaches}
\label{sec:approaches}

All text must be in two-column format.

Detail the initial approaches, architectures, and methods you experimented with when first tackling the problem.

Discuss what worked well during your experimentation and what did not, including:

\begin{itemize}
    \item \textbf{Feature Engineering \& Selection}: Explain the features you created, why you believed they
would be predictive, and the methods you used to select your final feature set.
    \item \textbf{Model Experimentation \& Refinement}: Discuss the variety of modeling approaches you
explored. Detail your process for optimizing your models and the techniques you used to
improve upon baseline performance, from simple models to potentially more complex
structures.
    \item \textbf{Validation Strategy}: Describe the validation strategy you used to make your choices and analyze its
effectiveness in predicting your final performance on the public and private leaderboards.

\end{itemize}

\subsection{References}

List and number all bibliographical references in 9-point Times, single-spaced, at the end of your paper.
When referenced in the text, include the citation number in square brackets, for
example~\cite{Authors14}.
Where appropriate, include page numbers and the name(s) of editors of referenced books.
When citing multiple papers at once, make sure that you cite them in numerical order such as this \cite{Alpher02,Alpher03,Alpher05,Authors14b,Authors14}.
If you use the template as advised, this will be taken care of automatically.

\begin{table}[h!]
  \caption{Results.   Ours is better.}
  \label{tab:example}
  \centering
  \begin{tabular}{@{}lc@{}}
    \toprule
    Method & Frobnability \\
    \midrule
    Theirs & Frumpy \\
    Yours & Frobbly \\
    Ours & Makes one's heart Frob\\
    \bottomrule
  \end{tabular}
\end{table}

\subsection{Illustrations, graphs, and photographs}

All graphics should be centered.
In \LaTeX, avoid using the \texttt{center} environment for this purpose, as this adds potentially unwanted whitespace.
Instead, use
{\small\begin{verbatim}
  \centering
\end{verbatim}}
at the beginning of your figure.
Please ensure that any point you wish to make is resolvable in a printed copy of the paper.
Resize fonts in figures to match the font in the body text, and choose line widths that render effectively in print.
Readers (and reviewers), even of an electronic copy, may choose to print your paper in order to read it.
You cannot insist that they do otherwise and, therefore, must not assume that they can zoom in to see tiny details on a graphic.

When placing figures in \LaTeX, it is almost always best to use \verb+\includegraphics+, and to specify the figure width as a multiple of the line width as in the example below
{\small\begin{verbatim}
   \usepackage{graphicx} ...
   \includegraphics[width=0.8\linewidth]
                   {myfile.pdf}
\end{verbatim}
}